\documentclass[11pt,a4paper]{article}
\usepackage[utf8]{inputenc}
\usepackage[T1]{fontenc}
\usepackage[brazil]{babel}
\usepackage[margin=2.5cm]{geometry}
\usepackage{booktabs}
\usepackage{amsmath,amssymb}
\usepackage[hidelinks]{hyperref}

\title{A Transformada Quântica de Fourier no teto de 28 qubits\\
\large Metodologia, verificação numérica e comparativo com o \texttt{cuStateVec}}
\author{RGPU --- simulação quântica em WGSL portátil (\texttt{rqubit})}
\date{13 de setembro de 2026}

\begin{document}
\maketitle

\begin{abstract}
Este relatório documenta a execução da Transformada Quântica de Fourier (QFT)
completa --- o circuito canônico de benchmark, com $O(n^2)$ portas --- no maior
estado que o binding do WebGPU permite: 27 qubits em precisão simples e
\textbf{28 qubits em meia precisão}, 420 portas, cada uma uma passada completa
por 1~GiB de estado. A carga foi medida no \texttt{rqubit} (kernels WGSL sobre
Vulkan, sem CUDA) e no \texttt{cuStateVec}, a biblioteca de vetor de estado da
própria NVIDIA, na mesma máquina. Resultados principais: empate técnico com o
\texttt{cuStateVec} em precisão simples (ambos a $\sim$224--228~GB/s), ganho de
$2{,}43\times$ em tempo e $2{,}44\times$ em energia com a fusão de portas até
6 qubits, e $1{,}96\times$ em tempo no teto de 28 qubits, onde a meia precisão
faz o papel que o \texttt{cuStateVec} não oferece. Dois resultados negativos
foram medidos e registrados: a fusão até 6 qubits \emph{perde} em 22 qubits, e
a fusão de 2 qubits em \texttt{f16} economiza tempo e \emph{gasta} energia.
\end{abstract}

\section{Introdução e objetivo}

Portas isoladas medem o kernel; um algoritmo mede a pilha. O objetivo deste
experimento é submeter o \texttt{rqubit} --- simulador de vetor de estado
quântico em Rust puro, sobre \texttt{wgpu}/WGSL --- a uma carga com a estrutura
de um algoritmo de verdade e compará-lo, na mesma máquina e na mesma carga, ao
\texttt{cuStateVec} da NVIDIA. A QFT foi escolhida por três razões:

\begin{itemize}
\item é o circuito de benchmark canônico, usado em demonstrações do Qiskit e do
  cuQuantum, o que torna o comparativo reconhecível;
\item tem fator de ramificação alto: fases controladas entre \emph{todos} os
  pares de qubits, o que estressa tanto o kernel de 2 qubits quanto a fusão;
  \item escala com $O(n^2)$ portas, o que permite chegar ao teto de qubits da
  máquina sem explosão de tempo.
\end{itemize}

O experimento também serve de primeiro uso real do produto interno em GPU
($\langle a|b\rangle$), implementado em dois estágios sobre o mesmo vetor de
estado: é ele que prova, depois da QFT completa, que a norma do estado sobreviveu
a 420 portas.

\section{Metodologia}

\subsection{O circuito}

A QFT segue a forma padrão com o qubit $q$ de peso $2^q$:
\begin{enumerate}
\item para cada $j \in \{0,\dots,n-1\}$: Hadamard em $q_j$, seguida de portas
  controladas-fase $\mathrm{CP}(2\pi/2^{k-j})$ entre todo par $(q_k, q_j)$ com
  $k > j$;
\item reversão final da ordem dos qubits por SWAPs.
\end{enumerate}
O total é $n + n(n-1)/2 + n/2$ portas --- 264 em 22 qubits, 364 em 26, 391 em
27 e \textbf{420 em 28}. A inversa exata, usada na verificação de ida-e-volta,
inverte a ordem das portas e troca o sinal dos ângulos.

A porta controlada-fase $\mathrm{CP}(\theta) = \operatorname{diag}(1,1,1,e^{i\theta})$
não existe pronta na API; foi construída diretamente sobre o campo público de
matriz de \texttt{Porta2}, com a fase no elemento $[3][3]$ --- a convenção de
índices ($|11\rangle$ por último) foi verificada no kernel e em
\texttt{aplicar\_cpu}.

\subsection{A pilha e as precisões}

O \texttt{rqubit} guarda o vetor de estado residente na GPU como um único
buffer. O teto de qubits é imposto pelo limite de binding do WebGPU (2~GB
menos 4 bytes):
\begin{itemize}
\item \texttt{f32} (\texttt{complex64}): 8~B por amplitude $\Rightarrow$
  \textbf{27 qubits} (1~GiB);
\item \texttt{f16} empacotado: 4~B por amplitude $\Rightarrow$ \textbf{28 qubits}
  (1~GiB) --- 29 qubits exigiriam exatamente 2~GiB, acima do limite.
\end{itemize}

Cada porta é uma passada completa pelo vetor de estado; numa carga limitada por
banda de memória, o tempo por porta é governado pelos bytes movidos
(leitura + escrita). A fusão de portas reduz o número de passadas: a fusão de
2 qubits absorve portas de 1 qubit pendentes na porta de 2 seguinte; a fusão
até 6 qubits usa os kernels especializados (amplitudes em registradores,
índices literais), que só existem em \texttt{f32}.

\subsection{Protocolo de medição}

\begin{itemize}
\item \textbf{Tempo:} uma execução de aquecimento, depois três execuções
  cronometradas; o valor reportado é a média por circuito. Cada execução
  aplica a QFT inteira sobre o estado corrente --- o custo das portas não
  depende do estado.
\item \textbf{Energia do \texttt{rqubit}:} janela calibrada de pelo menos
  $\sim$4~s (repetições ajustadas por circuito, mínimo de 3), medida pelo
  \texttt{rgpu-power} via \texttt{nvidia-smi}, com ociosidade calibrada antes
  de cada rodada.
\item \textbf{Energia do \texttt{cuStateVec}:} medida por fora do processo
  (\texttt{medir\_externo}), pela diferença entre uma execução só de partida
  (\texttt{reps=0}) e a execução com $R$ repetições, dividida pelo número de
  circuitos de fato executados (aquecimento incluso). Os dois instrumentos não
  são idênticos; os números de energia são comparáveis com ressalva.
\item \textbf{Comparativo:} \texttt{cuStateVec} via \texttt{cuquantum-python}
  (precisão simples, \texttt{complex64}), matrizes pré-carregadas no
  dispositivo, venv dedicado fora do repositório. O Aer (CPU) não entrou nesta
  rodada; a GPU dele não foi instalada (incompatível com o Qiskit 2.5) e a
  comparação GPU$\times$CPU já está coberta no README para portas isoladas.
\end{itemize}

\section{Testes de verificação}

A medição só vale se a física estiver de pé. Três verificações acompanham a
execução, todas com o produto interno em GPU como instrumento:

\begin{table}[h]
\centering
\begin{tabular}{llrl}
\toprule
Teste & O que prova & Medido & Tolerância \\
\midrule
8 qubits: QFT GPU contra CPU & correção do circuito & erro máx.\ $1{,}12\cdot10^{-8}$ & $10^{-4}$ \\
8 qubits: $|{\rm amplitude}|$ & estrutura da QFT de $|x_0\rangle$ & $0{,}0625$ exato & $1/\sqrt{N}$ \\
27 qubits \texttt{f32}: $\langle\psi|\psi\rangle$ após QFT & normalização & $0{,}999997$ & $10^{-4}$ \\
27 qubits \texttt{f32}: $\langle x_0|\mathrm{QFT}^{-1}\mathrm{QFT}|x_0\rangle$ & cadeia de 782 portas & $0{,}999998$ & $10^{-4}$ \\
28 qubits \texttt{f16}: $\langle\psi|\psi\rangle$ após QFT & normalização em \texttt{f16} & $0{,}997077$ & $10^{-2}$ \\
28 qubits \texttt{f16}: ida-e-volta (785 portas fundidas) & física em \texttt{f16} & $0{,}998535$ & $10^{-2}$ \\
\bottomrule
\end{tabular}
\caption{Verificações numéricas. O estado inicial $|x_0\rangle$ é preparado com
portas X (padrão de bits alternados; em 8 qubits, $x_0 = 0\mathrm{b}10110001$).}
\end{table}

O desvio do \texttt{f16} ($\approx -2{,}9\cdot10^{-3}$ na norma) é o erro de
arredondamento acumulado de centenas de portas cujos ângulos não são
representáveis em meia precisão --- mesma ordem de grandeza medida em testes
unitários de rotações. Um efeito estrutural vem de graça: ângulos CP abaixo da
resolução de \texttt{f16} ($\sim5\cdot10^{-4}$) agem como identidade, de modo
que a QFT em \texttt{f16} é \emph{aproximada por construção} --- a QFT
aproximada da literatura, sem custo extra de implementação.

\section{Resultados de tempo}

\begin{table}[h]
\centering
\begin{tabular}{rrrrrrr}
\toprule
Qubits & Prec. & Portas & Fundidas & ms/circuito & ms/porta & Banda (GB/s) \\
\midrule
22 & f32 & 264 & 264 (direto) & 67,6 & 0,256 & --- \\
22 & f32 & 264 & 242 (fusão-2) & 55,8 & 0,230 & --- \\
22 & f32 & 264 & 56 (fusão-6) & 90,7 & 1,620 & --- \\
24 & f32 & 312 & 312 (direto) & 392,2 & 1,257 & --- \\
24 & f32 & 312 & 60 (fusão-6) & \textbf{199,8} & 3,329 & --- \\
26 & f32 & 364 & 364 (direto) & 1.758,7 & 4,832 & 212 \\
26 & f32 & 364 & 68 (fusão-6) & \textbf{756,2} & 11,121 & 494 \\
26 & f16 & 364 & 364 (direto) & \textbf{925,2} & 2,542 & 403 \\
26 & f16 & 364 & 338 (fusão-2) & 822,4 & 2,433 & 455 \\
27 & f32 & 391 & 391 (direto) & 3.745,2 & 9,579 & 224 \\
27 & f32 & 391 & 82 (fusão-6) & \textbf{1.539,7} & 18,777 & 524 \\
28 & f16 & 420 & 420 (direto) & 4.031,6 & 9,599 & 224 \\
28 & f16 & 420 & 392 (fusão-2) & 3.763,9 & 9,602 & 240 \\
\bottomrule
\end{tabular}
\caption{QFT completa no \texttt{rqubit}. A ``banda'' é o tráfego efetivo
($2\times$ o estado por porta, $2\times$ as portas fundidas por circuito),
não um teto de hardware. Em \texttt{f16}, \texttt{28} qubits não têm
contraparte \texttt{f32} (o binding não permite), e a fusão-6 não existe
(os kernels de $N$ qubits são só em \texttt{f32}).}
\end{table}

O caso direto é governado por banda: em 27 e 28 qubits o tráfego efetivo fica
em $\sim$224~GB/s de $\sim$256 disponíveis. O previsto pela banda
(420 portas $\times$ 2~GiB $\div$ 200~GB/s $\approx 4{,}2$~s) deu
4,03~s medidos --- dentro de 5\%.

\section{Resultados de energia}

\begin{table}[h]
\centering
\begin{tabular}{rrrrr}
\toprule
Qubits & Prec. & rqubit direto & rqubit fusão-6 & cuStateVec (por fora) \\
\midrule
26 & f32 & 120,9 J & \textbf{47,8 J} & $\sim$65,5 J \\
26 & f16 & \textbf{36,6 J} & --- & $\sim$65,5 J \\
27 & f32 & 249,0 J & \textbf{102,0 J} & --- \\
28 & f16 & \textbf{246,4 J} & --- & $\sim$288,3 J \\
\bottomrule
\end{tabular}
\caption{Energia por circuito. A coluna do \texttt{cuStateVec} vem de
medição externa (diferença entre execução de partida e execução com carga,
dividida pelos circuitos executados); a do \texttt{rqubit}, do medidor
integrado. Métodos distintos: comparar com ressalva.}
\end{table}

Dois destaques. A meia precisão faz \textbf{3,3$\times$ menos energia} que a
simples no mesmo circuito (36,6 contra 120,9 J em 26 qubits) --- mais que os
$2\times$ dos bytes, porque a potência média também cai. E a fusão de 2 qubits
em \texttt{f16} \emph{custa} energia: 822~ms contra 925 (ganho de 11\%), mas
48,9~J contra 36,6~J (+34\%). A Hadamard absorvida no CP faz aritmética extra
por byte; num regime limitado por banda isso não compra tempo --- cobra watt.
É a quarta vez neste repositório em que a medição de energia contradiz a de
tempo.

\section{Comparativo com o cuStateVec}

\begin{table}[h]
\centering
\begin{tabular}{rrrrrr}
\toprule
Qubits & rqubit direto & rqubit fusão-2 & rqubit fusão-6 & cuStateVec & Melhor \\
\midrule
22 & 67,6 & 55,8 & 90,7 & \textbf{37,7} & cuStateVec $1{,}8\times$ \\
24 & 392,2 & 362,1 & \textbf{199,8} & 367,2 & rqubit $1{,}84\times$ \\
26 & 1.758,7 & 1.636,1 & \textbf{756,2} & 1.716,2 & rqubit $2{,}27\times$ \\
26 (f16) & 925,2 & 822,4 & --- & 1.716,2 & rqubit $2{,}09\times$ \\
27 & 3.745,2 & 3.480,5 & \textbf{1.539,7} & 3.685,0 & rqubit $2{,}39\times$ \\
28 (f16) & 4.031,6 & \textbf{3.763,9} & --- & 7.916,4 & rqubit $1{,}96\times$ \\
\bottomrule
\end{tabular}
\caption{ms por circuito (QFT completa). \texttt{cuStateVec} em precisão
simples; não oferece meia precisão para vetor de estado.}
\end{table}

\textbf{Leitura 1 --- empate técnico em precisão simples.} Sem fusão, o
\texttt{rqubit} mede 3.745~ms contra 3.685 do \texttt{cuStateVec} em 27
qubits ($1{,}6\%$), e 1.759 contra 1.716 em 26. Os dois saturam a banda de
DRAM ($\sim$224--228~GB/s calculados); quando o gargalo é mover bytes, um
kernel WGSL portátil e a biblioteca proprietária da NVIDIA chegam ao mesmo
lugar. É o mesmo resultado do README para portas isoladas, agora num
algoritmo completo.

\textbf{Leitura 2 --- a fusão é o diferencial.} Os CPs consecutivos sobre o
mesmo qubit alvo fundem-se em unitárias de até 6 qubits: 391 portas viram 82,
e o circuito cai $2{,}43\times$ em tempo e $2{,}44\times$ em energia. A
comparação é justa na direção do que cada API entrega: o \texttt{cuStateVec}
é de baixo nível e não funde portas; a fusão do \texttt{rqubit} é parte da
pilha medida. A NVIDIA tem camadas superiores (\texttt{cutensornet}) que
fariam otimização equivalente --- a ressalva já registrada no README vale
aqui.

\textbf{Leitura 3 --- o regime pequeno pertence ao cuStateVec.} Em 22 qubits o
estado (32~MiB) cabe na L2 da placa: as passadas ficam baratas, o custo fixo
por dispatch domina, e o \texttt{cuStateVec} --- com despacho mais leve
(0,143~ms/porta contra 0,256) --- vence $1{,}8\times$. Nesse regime, cortar
passadas não rende: a fusão até 6 qubits é 34\% \emph{mais lenta} que o
caminho direto. O cruzamento fica entre 22 e 24 qubits e não foi varrido.

\textbf{Leitura 4 --- meia precisão compra um qubit inteiro.} No teto, o
\texttt{rqubit} em \texttt{f16} faz a QFT de 28 qubits em 4,03~s contra 7,92~s
do \texttt{cuStateVec} em \texttt{f32} ($1{,}96\times$) e $\sim$288~J contra
246~J. A conta é de banda: ambos movem bytes a $\sim$224--228~GB/s; o
\texttt{f16} move metade dos bytes por porta e, portanto, cabe um estado com
o dobro de amplitudes no mesmo tempo.

\section{Ameaças à validade}

\begin{itemize}
\item \textbf{Deriva térmica:} GPU de notebook muda de clock sozinha; efeitos
  abaixo de $\sim$20\% não são resolvíveis por cronometragem de parede nesta
  máquina (registro em \texttt{RESULTADOS-NEGATIVOS.md}). As linhas do
  comparativo foram medidas em sessões separadas; o empate de 27 qubits
  ($1{,}6\%$) está dentro dessa margem e deve ser lido como empate, não como
  vitória do \texttt{cuStateVec}.
\item \textbf{Métodos de energia distintos:} medidor integrado
  (\texttt{rqubit}) contra medição externa do processo (\texttt{cuStateVec}),
  com janelas e composições de ociosidade diferentes. As diferenças de energia
  abaixo de $\sim$30\% não são conclusivas; as de $2\times$ ou mais, sim.
\item \textbf{Repetição sobre estado corrente:} as repetições aplicam a QFT
  sobre o estado já transformado. O custo das portas não depende do estado,
  mas o padrão de bits pode afetar consumo de forma residual.
\item \textbf{Água no aquecimento do cuStateVec:} a partida do Python e a
  inicialização do CUDA foram isoladas pela diferença entre \texttt{reps=0} e
  \texttt{reps}$>0$; o resíduo de partida é da ordem de 47~J.
\end{itemize}

\section{Conclusões}

No algoritmo canônico de benchmark, no maior estado que a máquina comporta:

\begin{enumerate}
\item o kernel WGSL portátil \textbf{empata com o \texttt{cuStateVec}} quando
  ambos limitados pela mesma banda de DRAM;
\item a \textbf{fusão de portas até 6 qubits} é o diferencial de arquitetura
  ($2{,}43\times$ tempo, $2{,}44\times$ energia acima de 24 qubits) --- e
  perde abaixo dele, onde a L2 muda o regime;
\item a \textbf{meia precisão} entrega um qubit a mais ($1{,}96\times$ contra
  o \texttt{cuStateVec}) e $3{,}3\times$ menos energia, ao custo de uma QFT
  aproximada cujo erro medido fica em $\sim10^{-3}$;
\item a \textbf{energia voltou a contradizer o tempo}: a fusão de 2 qubits em
  \texttt{f16} economiza passadas e gasta watts.
\end{enumerate}

\appendix
\section{Reprodutibilidade}

\begin{verbatim}
# rqubit: QFT completa (verificações + tabela; `tabela` refaz só as medições)
cargo run -p rqubit --release --example bench_qft
cargo run -p rqubit --release --example bench_qft -- tabela 27:f32 28:f16

# cuStateVec (venv dedicado, fora do repositório)
export LD_LIBRARY_PATH=$(find ~/.cache/rgpu-benchmarks/quantum \
  -name lib -maxdepth 2 -type d | tr '\n' ':')
python benchmarks/quantum/bench_custatevec_qft.py 28 3 single

# energia do processo externo
cargo run -p rqubit --release --example medir_externo -- \
  ~/.cache/rgpu-benchmarks/quantum/bin/python \
  benchmarks/quantum/bench_custatevec_qft.py 26 6 single
\end{verbatim}

Máquina: RTX 4070 Laptop (8~GB), Intel Meteor Lake, wgpu/Vulkan, driver NVIDIA
595.91.07. Versões: burn 0.21, TensorFlow 2.21.0, PyTorch 2.14.0, Qiskit 2.5,
\texttt{cuquantum-python-cu12}.

\end{document}
